\documentclass[10pt,a4paper]{article}
\usepackage[margin=1.8cm]{geometry}
\usepackage{amsmath,amssymb,amsthm}
\usepackage{physics}
\usepackage{enumitem}
\usepackage{xcolor}
\usepackage{tcolorbox}
\usepackage{booktabs}
\usepackage{multicol}
\usepackage{mdframed}
\usepackage{hyperref}
\usepackage{titlesec}
\usepackage{fancyhdr}
\usepackage{array}
\usepackage{tikz}
\usetikzlibrary{arrows.meta,calc}

% Kleuren
\definecolor{theoryblue}{RGB}{30,90,160}
\definecolor{exblue}{RGB}{0,120,190}
\definecolor{hintgreen}{RGB}{0,130,70}
\definecolor{warningred}{RGB}{180,30,30}
\definecolor{formulagray}{RGB}{245,245,250}
\definecolor{keybox}{RGB}{255,248,220}
\definecolor{headerbg}{RGB}{20,60,120}
\definecolor{tipsbox}{RGB}{230,255,230}

% tcolorbox stijlen
\tcbuselibrary{skins,breakable}

\newtcolorbox{theorievraag}[2][]{
  enhanced, breakable,
  title={#2},
  colback=blue!4!white,
  colframe=theoryblue,
  fonttitle=\bfseries\color{white},
  coltitle=white,
  attach boxed title to top left={yshift=-2mm,xshift=5mm},
  boxed title style={colback=theoryblue, rounded corners},
  #1
}

\newtcolorbox{oefening}[2][]{
  enhanced, breakable,
  title={#2},
  colback=green!4!white,
  colframe=hintgreen!80!black,
  fonttitle=\bfseries\color{white},
  coltitle=white,
  attach boxed title to top left={yshift=-2mm,xshift=5mm},
  boxed title style={colback=hintgreen!80!black, rounded corners},
  #1
}

\newtcolorbox{oplossing}[1][]{
  enhanced, breakable,
  title={Oplossing/Aanpak},
  colback=yellow!6!white,
  colframe=orange!70!black,
  fonttitle=\bfseries\small,
  coltitle=black,
  #1
}

\newtcolorbox{formulabox}[1][]{
  enhanced,
  colback=formulagray,
  colframe=gray!50,
  fonttitle=\bfseries\small,
  #1
}

\newtcolorbox{keyformulas}[1][]{
  enhanced,
  colback=keybox,
  colframe=orange!60!black,
  fonttitle=\bfseries,
  #1
}

\newtcolorbox{tip}[1][]{
  enhanced,
  colback=tipsbox,
  colframe=hintgreen,
  #1
}

% Headers
\pagestyle{fancy}
\fancyhf{}
\fancyhead[L]{\textcolor{headerbg}{\textbf{Thermische Fysica -- Examenvoorbereiding}}}
\fancyhead[R]{\textcolor{headerbg}{\thepage}}
\fancyfoot[C]{\textcolor{gray}{\small 2e Bachelor Fysica \& Sterrenkunde}}
\renewcommand{\headrulewidth}{0.4pt}
\renewcommand{\headrule}{\hbox to\headwidth{\color{headerbg}\leaders\hrule height\headrulewidth\hfill}}

% Sectie opmaak
\titleformat{\section}{\large\bfseries\color{headerbg}}{}{0em}{\colorbox{headerbg!15}{\parbox{\dimexpr\linewidth-2\fboxsep}{\textcolor{headerbg}{\thesection.\ #1}}}}
\titleformat{\subsection}{\normalsize\bfseries\color{theoryblue}}{}{0em}{\thesubsection.\ #1}

\setlist[itemize]{leftmargin=*,topsep=2pt,itemsep=1pt}
\setlist[enumerate]{leftmargin=*,topsep=2pt,itemsep=1pt}

% Nuttige macros
\newcommand{\dbar}{\bar{d}}
\newcommand{\pder}[2]{\left(\frac{\partial #1}{\partial #2}\right)}

\begin{document}

%===========================================================================
\begin{center}
{\Large\bfseries\color{headerbg} Examenvoorbereiding Thermische Fysica}\\[2pt]
{\large 2e Bachelor Fysica \& Sterrenkunde}\\[4pt]
{\normalsize\color{gray} Theorie (gesloten boek) $\cdot$ Oefeningen (open boek) $\cdot$ Geordend naar frequentie}
\end{center}
\vspace{-3mm}
\noindent\textcolor{headerbg}{\rule{\linewidth}{1.5pt}}

%===========================================================================
\section{Formubleblad: Sleuteluitdrukkingen}
%===========================================================================

\begin{multicols}{2}

\begin{keyformulas}[title=Eerste Wet]
$dU = \dbar Q + \dbar W$\quad (infinitesimaal)\\
$\dbar W = -P\,dV$\; (hydrostatisch)\\
$U_f - U_i = Q + W$\; (eindig)
\end{keyformulas}

\begin{keyformulas}[title=Ideaal Gas]
$PV = nRT$,\quad $R = 8{,}314$\,J/(mol\,K)\\
$\left(\frac{\partial U}{\partial V}\right)_T = 0$,\quad $U = U(T)$ \\
$C_P - C_V = nR$,\quad $\gamma = C_P/C_V$\\
Adiabaat: $PV^\gamma = \text{const}$,\; $TV^{\gamma-1}=\text{const}$
\end{keyformulas}

\begin{keyformulas}[title=Entropie \& 2e Wet]
$dS = \frac{\dbar Q_R}{T}$\; (reversibel)\\
$\sum \Delta S_{\text{univ}} \geq 0$\; (irreversibel: $>0$)\\
$S = k_B \ln\Omega$ \; (Boltzmann)\\
Carnot: $\eta_\text{Carnot} = 1 - T_K/T_H$
\end{keyformulas}

\begin{keyformulas}[title=TdS-vergelijkingen]
$TdS = C_V\,dT + T\pder{P}{T}_V dV$\\
$TdS = C_P\,dT - T\pder{V}{T}_P dP$
\end{keyformulas}

\begin{keyformulas}[title=Maxwell-betrekkingen]
$\pder{T}{V}_S = -\pder{P}{S}_V$ \quad
$\pder{T}{P}_S = \pder{V}{S}_P$\\
$\pder{S}{V}_T = \pder{P}{T}_V$ \quad
$\pder{S}{P}_T = -\pder{V}{T}_P$
\end{keyformulas}

\begin{keyformulas}[title=Energievergelijkingen]
$\pder{U}{V}_T = T\pder{P}{T}_V - P = \frac{T\beta}{\kappa} - P$\\
$C_P - C_V = \frac{TV\beta^2}{\kappa}$\\
$\beta = \frac{1}{V}\pder{V}{T}_P$,\quad $\kappa = -\frac{1}{V}\pder{V}{P}_T$
\end{keyformulas}

\begin{keyformulas}[title=Thermodynamische Functies]
$H = U + PV$,\; $dH = \dbar Q + V\,dP$,\; $\pder{H}{T}_P = C_P$\\
$F = U - TS$,\; $dF = -S\,dT - P\,dV$\\
$G = H - TS$,\; $dG = -S\,dT + V\,dP$\\
Joule-Kelvin: $\mu = \pder{T}{P}_H = \frac{1}{C_P}\!\left[T\pder{V}{T}_P - V\right]$
\end{keyformulas}

\begin{keyformulas}[title=Statistische Mechanica]
Boltzmann: $N_i = \frac{N}{Z}\,g_i e^{-\varepsilon_i/kT}$,\; $Z = \sum_i g_i e^{-\varepsilon_i/kT}$\\
$U = NkT^2\pder{\ln Z}{T}_V$,\quad $P = NkT\pder{\ln Z}{V}_T$\\
Ideaal monoatomisch: $Z = V\!\left(\frac{2\pi mkT}{h^2}\right)^{3/2}$\\
Equipartitie: $\langle\varepsilon\rangle = \frac{f}{2}kT$\; ($f$ kwadr.\ termen)
\end{keyformulas}

\begin{keyformulas}[title=Fermi-Dirac / Bose-Einstein]
FD: $N_i = \frac{g_i}{e^{(\varepsilon_i-\varepsilon_F)/kT}+1}$\\
BE: $N_i = \frac{g_i}{Ae^{\varepsilon_i/kT}-1}$\; (fotonen: $A=1$)\\
Planck: $w_\nu = \frac{8\pi h\nu^3/c^3}{e^{h\nu/kT}-1}$\\
Stefan-Boltzmann: $R = \sigma T^4$,\; $\sigma = \frac{2\pi^5 k^4}{15h^3c^2}$
\end{keyformulas}

\begin{keyformulas}[title=Transport]
Fick: $d = -D\,\frac{\partial n}{\partial x}$,\quad $\frac{\partial n}{\partial t} = D\,\frac{\partial^2 n}{\partial x^2}$\\
Fourier: $e = -\lambda\,\frac{\partial T}{\partial x}$,\quad $\frac{\partial T}{\partial t} = a\,\frac{\partial^2 T}{\partial x^2}$\\
Vrije weglengte: $\ell = \frac{1}{\sqrt{2}\,n\sigma}$\\
Maxwell: $\frac{dN_v}{dv} = \frac{2N}{\sqrt{2\pi}}\!\left(\frac{m}{kT}\right)^{3/2}\!v^2 e^{-mv^2/2kT}$
\end{keyformulas}

\begin{keyformulas}[title=Clapeyron / Faseover-gangen]
$\frac{dP}{dT} = \frac{L}{T(V_f-V_i)}$\; (Clapeyron)\\
Clapeyron-Clausius (damp): $\frac{dP}{dT} \approx \frac{LP}{RT^2}$\\
1e orde: $L = T\Delta S$,\; $G_i = G_f$
\end{keyformulas}

\end{multicols}

\vspace{2mm}
\noindent\textcolor{headerbg}{\rule{\linewidth}{1pt}}

%===========================================================================
\section{Theorie (Gesloten Boek) -- Geordend naar Frequentie}
%===========================================================================

%-------------------------------------------------------------------
\subsection{Caratheodory's Axioma en de Tweede Wet \textnormal{\color{warningred}[$\star\star\star\star$\ -- 4 examens]}}

\begin{theorievraag}{Formuleer het axioma van Caratheodory. Bewijs de stelling en leg het belang van de tweede wet uit. \textnormal{(2010, 2015, 2018, 2019)}}

\textbf{Axioma van Caratheodory:}
In de onmiddellijke nabijheid van om het even welke evenwichtstoestand van een systeem bevinden zich toestanden die \emph{niet bereikbaar zijn} via reversibele adiabatische processen.

\end{theorievraag}

\begin{oplossing}
\textbf{Stap 1: Stelsel met 2 onafhankelijke variabelen (eenvoudig geval).}\\
Voor een systeem met coördinaten $(T_e, X)$ geldt de eerste wet:
$$\dbar Q = \pder{U}{T_e}_X dT_e + \left[-Y + \pder{U}{X}_{T_e}\right]dX$$
Dit is een vergelijking in 2 variabelen. Een dergelijke pfaffiaanse differentiaalvorm heeft \emph{altijd} een integrerende factor, onafhankelijk van de 2e wet. De reversibele adiabaten $\sigma(T_e, X) = \text{const}$ bestaan dus sowieso.

\textbf{Stap 2: Stelsel met 3+ onafhankelijke variabelen (bijv.\ $(U, X, X')$).}\\
Veronderstel dat vanuit toestand $i$ \emph{twee} punten $f_1$ en $f_2$ (verticaal boven elkaar in het $UXX'$-diagram) bereikbaar zijn via reversibele adiabaten.\\
Dan: $i \to f_1 \to f_2 \to i$ vormt een cyclus waarbij:
\begin{itemize}
  \item $f_1 \to f_2$ is isochoor ($dX = dX' = 0$) $\Rightarrow$ geen arbeid, wel warmte-opname $Q > 0$
  \item $i \to f_1$ en $i \to f_2$: adiabatische processen, enkel arbeid
\end{itemize}
Resultaat: een cyclus waarbij $Q > 0$ volledig in arbeid omgezet wordt. \textbf{Dit is strijdig met de Kelvin-Planck formulering van de 2e wet.}\\
$\Rightarrow$ Op elke verticale slechts 1 punt bereikbaar via RA-processen. Deze punten liggen op een \textbf{RA-oppervlak} $\sigma(T_e, X, X') = \text{const}$.

\textbf{Stap 3: Integreerbaarheid van $\dbar Q$.}\\
Uit het bestaan van RA-oppervlakken volgt:
$$\dbar Q = \lambda(T_e)\, f(\sigma)\, d\sigma \quad \Rightarrow \quad \frac{\dbar Q}{\lambda(T_e)} = \text{totale differentiaal } dS$$
De integrerende factor $1/\lambda(T_e) = 1/(kT)$ is \emph{universeel} (geldt voor alle systemen) en hangt enkel van $T$ af.

\textbf{Belang 2e wet:} Voor slechts 2 variabelen bestaat de integrerende factor altijd (wiskunde). Voor $\geq 3$ variabelen is de 2e wet \emph{essentieel} om het bestaan van de integrerende factor te garanderen.
\end{oplossing}

%-------------------------------------------------------------------
\subsection{Entropie: Definitie, Berekening, Link met 2e Wet \textnormal{\color{warningred}[$\star\star\star\star$\ -- 4 examens]}}

\begin{theorievraag}{Definieer entropie vanuit twee standpunten. Bereken de entropie van een ideaal gas. Bespreek de link met de tweede wet. \textnormal{(2016, 2018, 2019, 2023)}}
\end{theorievraag}

\begin{oplossing}
\textbf{Macroscopische definitie (thermodynamisch):}
$$dS = \frac{\dbar Q_R}{T} \quad \Rightarrow \quad S_f - S_i = \int_i^{f,\text{reversibel}} \frac{\dbar Q}{T}$$
$S$ is een toestandsfunctie (pad-onafhankelijk voor reversibele processen). 
$\oint_R \frac{\dbar Q}{T} = 0$ (Clausius-theorema).

\textbf{Statistische definitie (microscopisch):}
$$S = k_B \ln\Omega$$
waarbij $\Omega$ = het aantal microtoestanden (thermodynamische waarschijnlijkheid).
Entropie = maat voor \emph{wanorde}. Maximum bij thermisch evenwicht.

\textbf{Entropie van een ideaal gas:}\\
Gebruik de eerste wet: $\dbar Q = C_V\,dT + P\,dV$, deel door $T$:
$$dS = \frac{C_V}{T}dT + \frac{nR}{V}dV$$
Integreer:
$$\boxed{S = C_V \ln T + nR \ln V + S_0}$$
Of equivalent (Sackur-Tetrode formule uit statistische mechanica):
$$S = Nk\left[\frac{3}{2}\ln T + \ln\frac{V}{N} + \ln\!\left(\frac{2\pi mk}{h^2}\right)^{3/2} + \frac{5}{2}\right]$$

\textbf{Link met de 2e wet:}
\begin{itemize}
  \item Reversibel proces: $\Delta S_\text{univ} = 0$
  \item Irreversibel (spontaan) proces: $\Delta S_\text{univ} > 0$
  \item Entropie-principe: $\sum \Delta S \geq 0$ (universum neemt toe of blijft gelijk)
\end{itemize}
Dit is equivalent met Kelvin-Planck en Clausius-formulering.
\end{oplossing}

%-------------------------------------------------------------------
\subsection{Carnotcyclus, Kelvintemperatuurschaal en Rendement \textnormal{\color{warningred}[$\star\star\star\star$\ -- 4+ examens]}}

\begin{theorievraag}{Definieer de Kelvin-temperatuurschaal. Toon aan dat deze gelijk is aan de ideaalgastemperatuur. Bespreek de Carnotcyclus en het maximale rendement. \textnormal{(2009, 2011, 2018, 2019, 2022)}}
\end{theorievraag}

\begin{oplossing}
\textbf{Kelvin-temperatuurschaal:}\\
Vertrek van de verhouding van warmtehoeveelheden bij twee reversibele isothermen die dezelfde adiabatische oppervlakken verbinden:
$$\frac{Q}{Q'} = \frac{\phi(T_e)}{\phi(T'_e)}$$
Definieer $T_k / T'_k = Q/Q'$, met referentie $T'_k = 273{,}16\,\text{K}$ (tripelpunt water).

\textbf{Gelijkheid met ideaalgastemperatuur:}\\
Voor een ideaal gas: $Q = nRT\ln(V_c/V_b)$ en $Q' = nRT'\ln(V_d/V_a)$.\\
Via de adiabaten $TV^{\gamma-1}=\text{const}$ volgt $\ln(V_c/V_b) = \ln(V_d/V_a)$, dus $Q/Q' = T/T'$.\\
$\Rightarrow T_k = T_{\text{ideaalgas}}$.

\textbf{Carnotcyclus (reversibele machine tussen 2 reservoirs):}
\begin{enumerate}
  \item Isotherme expansie bij $T_H$: warmte-opname $|Q_H|$
  \item Adiabatische expansie: $T_H \to T_K$
  \item Isotherme compressie bij $T_K$: warmte-afgifte $|Q_K|$
  \item Adiabatische compressie: $T_K \to T_H$
\end{enumerate}
$$\eta_\text{Carnot} = 1 - \frac{|Q_K|}{|Q_H|} = 1 - \frac{T_K}{T_H}$$
\textbf{Dit is het maximale rendement van alle machines tussen $T_H$ en $T_K$.}\\
Bewijs: via entropie-principe: $\eta \leq 1 - T_K/T_H$ (uit $\Delta S_\text{univ} \geq 0$).\\
$\eta = 1$ is onmogelijk want dan $T_K = 0$ (3e wet: onbereikbaar).

\textbf{TS-diagram van Carnotcyclus:} rechthoek met zijden $T_H, T_K$ en $S_1, S_2$.\\
$|W| = (T_H - T_K)(S_2 - S_1)$,\; $|Q_H| = T_H(S_2-S_1)$.
\end{oplossing}

%-------------------------------------------------------------------
\subsection{Twee Formuleringen van de 2e Wet en Equivalentie \textnormal{\color{warningred}[$\star\star\star$\ -- 3 examens]}}

\begin{theorievraag}{Geef de formuleringen van Kelvin-Planck en Clausius. Bewijs dat ze equivalent zijn. \textnormal{(2009, 2018, 2019)}}
\end{theorievraag}

\begin{oplossing}
\textbf{Kelvin-Planck (K):} Er bestaat geen enkel proces dat als enig resultaat heeft: warmte onttrekken aan één reservoir en dit volledig omzetten in arbeid.\\
\textbf{Clausius (C):} Er bestaat geen enkel proces dat als enig resultaat heeft: warmte overbrengen van een koud naar een warm reservoir zonder externe arbeid.

\textbf{Bewijs $\neg K \Rightarrow \neg C$:}\\
Stel K is geldig, dus een machine $M_1$ onttrekt $|Q_1|$ aan warme bron en levert $|W| = |Q_1|$.\\
Gebruik $|W|$ voor een warmtepomp $M_2$: $|Q_2|$ uit koude bron, $|Q_1|+|Q_2|$ naar warme bron.\\
Netto: $|Q_2|$ van koud naar warm zonder arbeid. $\Rightarrow \neg C$. \checkmark

\textbf{Bewijs $\neg C \Rightarrow \neg K$:}\\
Stel C is geldig: $|Q_2|$ van koud naar warm zonder arbeid.\\
Laat machine $M$: $|Q_1|$ van warm, $|Q_2|$ naar koud, arbeid $|W| = |Q_1|-|Q_2|$.\\
Netto: $|Q_1|-|Q_2|$ onttrokken aan warm, volledig in arbeid. $\Rightarrow \neg K$. \checkmark
\end{oplossing}

%-------------------------------------------------------------------
\subsection{TdS-Vergelijkingen \textnormal{\color{warningred}[$\star\star\star$\ -- 3 examens]}}

\begin{theorievraag}{Stel de TdS-vergelijkingen op voor een hydrostatisch systeem. Hoe zien ze eruit voor een ideaal gas? Wat zijn de toepassingen? \textnormal{(2009, 2023, 2024 2e zit)}}
\end{theorievraag}

\begin{oplossing}
\textbf{Eerste TdS-vergelijking} ($S$ als functie van $T$ en $V$):
$$\boxed{TdS = C_V\,dT + T\pder{P}{T}_V dV}$$
Afleiden: $dS = \pder{S}{T}_V dT + \pder{S}{V}_T dV$. Gebruik $C_V = T\pder{S}{T}_V$ en de Maxwell-betrekking $\pder{S}{V}_T = \pder{P}{T}_V$.

\textbf{Tweede TdS-vergelijking} ($S$ als functie van $T$ en $P$):
$$\boxed{TdS = C_P\,dT - T\pder{V}{T}_P dP}$$
Gebruik $\pder{S}{P}_T = -\pder{V}{T}_P$ (Maxwell).

\textbf{Voor een ideaal gas} ($PV = nRT$): $\pder{P}{T}_V = nR/V = P/T$ en $\pder{V}{T}_P = nR/P = V/T$:
$$TdS = C_V\,dT + nR\,\frac{dV}{V} \cdot T \quad \Rightarrow \quad dS = \frac{C_V}{T}dT + \frac{nR}{V}dV$$
$$TdS = C_P\,dT - nR\,\frac{dP}{P} \cdot T$$

\textbf{Toepassingen:}
\begin{itemize}
  \item Joule-Kelvin coëfficiënt: uit $dH = TdS + VdP$ en 2e TdS-vergelijking: $\mu = \frac{1}{C_P}\left[T\pder{V}{T}_P - V\right]$
  \item Adiabatische drukverandering: $\Delta T = \frac{TV\beta}{C_P}\Delta P$
  \item Isotherme drukverandering: $Q = -TV\beta\,\Delta P$
\end{itemize}
\end{oplossing}

%-------------------------------------------------------------------
\subsection{Stralingswet van Planck en Stefan-Boltzmann \textnormal{\color{warningred}[$\star\star\star$\ -- 3 examens]}}

\begin{theorievraag}{Bespreek de stralingswet van Planck. Leid de constante van Stefan-Boltzmann af. \textnormal{(2010, 2011, 2018, 2019)}}
\end{theorievraag}

\begin{oplossing}
\textbf{Fotonengas en Bose-Einstein statistiek:}\\
Fotonen zijn bosonen, hun aantal is niet geconserveerd ($A=1$). Bezetting:
$$N_\nu = \frac{g_\nu}{e^{h\nu/kT}-1}$$
Ontaarding (staande golven in volume $V$, 2 polarisaties):
$$g_\nu = \frac{8\pi\nu^2 V}{c^3}$$

\textbf{Energiedichtheid (Planck):}
$$\boxed{w_\nu = \frac{8\pi h\nu^3/c^3}{e^{h\nu/kT}-1}}$$

\textbf{Stefan-Boltzmann via thermodynamica:}\\
Fotonengas: $P = w/3$ (stralingsdruk), $U = wV$.\\
Energievergelijking: $\pder{U}{V}_T = T\pder{P}{T}_V - P$
$$w = \frac{T}{3}\frac{dw}{dT} - \frac{w}{3} \quad \Rightarrow \quad \frac{dw}{w} = 4\frac{dT}{T} \quad \Rightarrow \quad w = bT^4$$
Met $R = cw/4$: $R = \sigma T^4$ met $\boxed{\sigma = \frac{2\pi^5 k^4}{15h^3c^2} = 5{,}67 \times 10^{-8}\,\text{W\,m}^{-2}\text{K}^{-4}}$.

\textbf{Wet van Wien:} maximum van $w_\lambda$: $\lambda_\text{max}T = 2{,}898 \times 10^{-3}\,\text{m\,K}$.
\end{oplossing}

%-------------------------------------------------------------------
\subsection{Partitiefunctie en Equipartitiewet \textnormal{\color{warningred}[$\star\star\star$\ -- 3 examens]}}

\begin{theorievraag}{Leid de partitiefunctie af voor een ideaal monoatomisch gas. Bewijs de Sackur-Tetrode formule. Leid de equipartitiewet af. \textnormal{(2011, 2018, 2019)}}
\end{theorievraag}

\begin{oplossing}
\textbf{Partitiefunctie:} $Z = \sum_\text{toestanden} e^{-\varepsilon_j/kT}$ met $\varepsilon_j = \frac{h^2}{8mV^{2/3}}(n_x^2+n_y^2+n_z^2)$.\\
Vervang som door integraal (niveaus liggen dicht):
$$Z = \int_0^\infty e^{-\frac{h^2 n_x^2}{8mkTa^2}} dn_x \cdots = V\!\left(\frac{2\pi mkT}{h^2}\right)^{3/2}$$

\textbf{Thermodynamische grootheden:}
$$U = NkT^2\pder{\ln Z}{T}_V = \frac{3}{2}NkT, \quad P = NkT\pder{\ln Z}{V}_T = \frac{NkT}{V}$$
$$S = Nk\ln\frac{Z}{N} + \frac{U}{T} + Nk = Nk\!\left[\frac{3}{2}\ln T + \ln\frac{V}{N} + \ln\!\left(\frac{2\pi mk}{h^2}\right)^{3/2} + \frac{5}{2}\right]$$
Dit is de \textbf{Sackur-Tetrode formule}.

\textbf{Equipartitiewet:} Als $\varepsilon = \sum_{i=1}^f b_i p_i^2$ (som van $f$ kwadratische termen):
$$Z = \prod_i \int_0^\infty e^{-\beta b_i p_i^2} dp_i = \beta^{-f/2} \prod_i K_i$$
$$\langle\varepsilon\rangle = -\frac{\partial \ln Z}{\partial\beta} = \frac{f}{2\beta} = \frac{f}{2}kT$$
\textbf{Elke kwadratische vrijheidsgraad draagt $\frac{1}{2}kT$ bij.}
\end{oplossing}

%-------------------------------------------------------------------
\subsection{Bose-Einstein Statistiek \textnormal{\color{warningred}[$\star\star$\ -- 2 examens]}}

\begin{theorievraag}{Leg de Bose-Einstein statistiek uit en teken de verdelingsfunctie. Vergelijk met Fermi-Dirac. \textnormal{(2009, 2018, 2019)}}
\end{theorievraag}

\begin{oplossing}
\textbf{Bosonen:} Deeltjes met heeltallige spin (He-4, H$_2$, fotonen, mesonen). Geen uitsluitingsprincipe: meerdere deeltjes in dezelfde toestand.

\textbf{Telmethode:} $N_i$ deeltjes over $g_i$ cellen:
$(N_i + g_i - 1)!$ rangschikkingen, deel door $N_i!(g_i-1)!$:
$$\Omega_{BE} = \prod_i \frac{(N_i+g_i-1)!}{N_i!(g_i-1)!}$$
Maximaliseren $\Rightarrow$:
$$\boxed{N_i = \frac{g_i}{Ae^{\varepsilon_i/kT}-1}}$$

\textbf{Fermi-Dirac} (fermionen: elektronengetal+spin 1/2, uitsluitingsprincipe $N_i \leq g_i$):
$$N_i = \frac{g_i}{e^{(\varepsilon_i-\varepsilon_F)/kT}+1}$$
Bij $T=0$: alle niveaus tot $\varepsilon_F$ bezet. $\varepsilon_F \sim 5\,\text{eV}$ voor metalen $\Rightarrow \Theta_F \sim 60\,000\,\text{K}$.

\textbf{Boltzmann-limiet:} $A \gg 1$ (lage dichtheid, hoge $T$): $N_i \approx \frac{g_i}{A}e^{-\varepsilon_i/kT}$ (beide statistieken geven hetzelfde).
\end{oplossing}

%-------------------------------------------------------------------
\subsection{Ideaalgas- en Kelvintemperatuur: Definitie en Gelijkheid \textnormal{\color{warningred}[$\star\star$\ -- 2 examens]}}

\begin{theorievraag}{Definieer de ideaalgastemperatuur en de Kelvintemperatuur. Toon aan dat ze gelijk zijn. \textnormal{(2018, 2019 mondeling)}}
\end{theorievraag}

\begin{oplossing}
\textbf{Ideaalgastemperatuur:} via constant-volume gasthermometer:
$$T = 273{,}16\,\text{K} \cdot \lim_{P_{TP}\to 0}\frac{P}{P_{TP}}$$

\textbf{Kelvintemperatuur:} Twee isothermen verbinden dezelfde twee RA-oppervlakken:
$$\frac{T_k}{T'_k} = \frac{Q}{Q'}$$
Referentie: $T'_k = 273{,}16\,\text{K}$ (tripelpunt).

\textbf{Bewijs gelijkheid:} Voor een ideaal gas volgen beide adiabaten $TV^{\gamma-1} = \text{const}$, zodat $\ln(V_c/V_b) = \ln(V_d/V_a)$ en dus $Q/Q' = T/T'$. De twee schalen vallen samen.
\end{oplossing}

%-------------------------------------------------------------------
\subsection{Methode van Rüchhardt \textnormal{\color{warningred}[$\star\star$\ -- 2 examens]}}

\begin{theorievraag}{Bespreek de methode van Rüchhardt voor de bepaling van $\gamma$. \textnormal{(2018, 2019 schriftelijk)}}
\end{theorievraag}

\begin{oplossing}
\textbf{Principe:} kogel (massa $m$, doorsnede $A$) in buis boven vat (volume $V$, druk $P = P_\text{atm} + mg/A$).\\
Bij kleine verplaatsing $y$ (adiabatisch, quasistatisch): $PV^\gamma = \text{const}$
$$P\gamma V^{\gamma-1}\Delta V + V^\gamma\Delta P = 0 \quad \Rightarrow \quad \Delta P = -\frac{\gamma P}{V}\Delta V = -\frac{\gamma P A^2}{V}y$$
Kracht: $F = A\Delta P = -\frac{\gamma P A^2}{V}y$ (terugdrijvend). Harmonische oscillatie met:
$$\tau = 2\pi\sqrt{\frac{mV}{\gamma P A^2}} \quad \Rightarrow \quad \boxed{\gamma = \frac{4\pi^2 mV}{\tau^2 A^2 P}}$$
Meet periode $\tau$, bereken $\gamma$.
\end{oplossing}

%-------------------------------------------------------------------
\subsection{Thermische Arbeidsmachines: Stirling en Benzinemotor \textnormal{\color{warningred}[$\star\star$\ -- 2 examens]}}

\begin{theorievraag}{Bespreek de Stirlingmotor/warmtepomp. Bespreek de benzinemotor (Ottocyclus). \textnormal{(2018, 2019 mondeling)}}
\end{theorievraag}

\begin{oplossing}
\textbf{Stirlingmotor (4 processen):}
\begin{enumerate}
  \item Isotherme compressie bij $T_K$: $|Q_K|$ afgifte
  \item Isochore verwarming via regenerator: opname $|Q_R|$
  \item Isotherme expansie bij $T_H$: $|Q_H|$ opname
  \item Isochore afkoeling via regenerator: afgifte $|Q_R|$
\end{enumerate}
Rendement (geïdealiseerd): $\eta = 1 - T_K/T_H$ = Carnot-rendement!\\
Regenerator = thermische buffer die $|Q_R|$ opslaat en teruggeeft.

\textbf{Ottocyclus (benzinemotor, lucht-gestandardiseerd):}
\begin{itemize}
  \item Adiabatische compressie $1\to2$: $T_1V_1^{\gamma-1} = T_2V_2^{\gamma-1}$
  \item Isochore warmteopname $2\to3$: $Q_H = C_V(T_3-T_2)$
  \item Adiabatische expansie $3\to4$
  \item Isochore warmteafgifte $4\to1$: $Q_K = C_V(T_1-T_4)$
\end{itemize}
$$\eta_\text{Otto} = 1 - \frac{T_4-T_1}{T_3-T_2} = 1 - \frac{1}{r^{\gamma-1}}, \quad r = V_1/V_2 \text{ (compressieverhouding)}$$
$r \approx 9$, $\gamma=1{,}4$: $\eta \approx 58\%$; reëel: $\approx 30\%$.
\end{oplossing}

%-------------------------------------------------------------------
\subsection{Maxwell Snelheidsverdeling \textnormal{\color{warningred}[$\star\star$\ -- 2 examens]}}

\begin{theorievraag}{Leid de Maxwell snelheidsverdeling af en definieer de karakteristieke snelheden. \textnormal{(2018, 2019 mondeling)}}
\end{theorievraag}

\begin{oplossing}
\textbf{Vertrek van energiedistributie:} $dN_\varepsilon = N\frac{dg_\varepsilon}{Z}e^{-\varepsilon/kT}$ met $dg_\varepsilon = \frac{2\pi V(2m)^{3/2}}{h^3}\varepsilon^{1/2}d\varepsilon$ en $Z = V\left(\frac{2\pi mkT}{h^2}\right)^{3/2}$.\\
Substitueer $\varepsilon = \frac{1}{2}mv^2$, $d\varepsilon = mv\,dv$:
$$\boxed{\frac{dN_v}{dv} = \frac{2N}{\sqrt{\pi}}\left(\frac{m}{2kT}\right)^{3/2}v^2 e^{-mv^2/2kT}}$$

\textbf{Karakteristieke snelheden:}
\begin{itemize}
  \item Meest waarschijnlijk: $v_p = \sqrt{2kT/m}$
  \item Gemiddeld: $\bar{v} = \sqrt{8kT/\pi m} \approx 1{,}13\,v_p$
  \item RMS: $v_\text{rms} = \sqrt{3kT/m} \approx 1{,}22\,v_p$
\end{itemize}
Verificatie met Estermann-Simpson-Stern experiment (cesium bundel, parabolische val).
\end{oplossing}

%-------------------------------------------------------------------
\subsection{Warmtegeleiding als Energietransport (Microscopisch) \textnormal{\color{warningred}[$\star\star$\ -- 2 examens]}}

\begin{theorievraag}{Bespreek warmtegeleiding als energietransport. Wat leren microscopische beschouwingen? \textnormal{(2022, 2023)}}
\end{theorievraag}

\begin{oplossing}
\textbf{Macroscopisch (wet van Fourier):} $e = -\lambda\pder{T}{x}$, diffusievergelijking: $\partial T/\partial t = a\,\partial^2T/\partial x^2$ met $a = \lambda/(\rho c_P)$.

\textbf{Microscopisch (voor gassen):}\\
Moleculen die vanuit vlak $x_0-\ell$ komen dragen energie $u(x_0-\ell)$:
$$e = \frac{1}{6}nv[u(x_0-\ell) - u(x_0+\ell)] = -\frac{1}{3}nv\ell\frac{\partial u}{\partial T}\frac{\partial T}{\partial x}$$
$$\lambda = \frac{1}{3}nv\ell c'_V = \frac{1}{3}n\bar{v}\ell\frac{c_V}{N_A}$$
Met $\ell = 1/(\sqrt{2}n\sigma)$: $n$ en $\ell$ heffen elkaar op $\Rightarrow$ \textbf{$\lambda$ is druk-onafhankelijk}.\\
$\lambda \propto \sqrt{T/m}/d^2$ met $d$ de molecuuldiameter.

\textbf{Vaste stoffen:} Via fononen (roostervibraties). Metalen: ook elektronen ($\lambda_e \propto T$).
\end{oplossing}

%-------------------------------------------------------------------
\subsection{Brownse Beweging en Diffusie van Einstein \textnormal{\color{warningred}[$\star\star$\ -- 1-2 examens]}}

\begin{theorievraag}{Leg de Brownse beweging uit. Leid de relatie $\langle l_B^2\rangle = 6\mu kT\,t$ af. \textnormal{(2018, 2019 mondeling)}}
\end{theorievraag}

\begin{oplossing}
\textbf{Fenomeen:} Kleine deeltjes ($\sim 1\,\mu$m) in vloeistof voeren zigzagbewegingen uit door ongelijke botsingen van vloeistofmoleculen. $\langle\varepsilon_\text{kin}\rangle = \frac{3}{2}kT$.

\textbf{Bewegingsvergelijking:} $m\ddot{x} = F'_x - \frac{1}{\mu}\dot{x}$ (viskeuze kracht + willekeurige kracht).\\
Na tijdsmiddeling ($\langle xF'_x\rangle = 0$, evenwicht):
$$m\frac{ds}{dt} + \frac{s}{\mu} = 2kT, \quad s = \langle d(x^2)/dt\rangle$$
Oplossing voor $t \gg \mu m$: $s = 2\mu kT$ $\Rightarrow$ $\langle x^2\rangle = 2\mu kT\,t$.

\textbf{In 3D:} $\langle l_B^2\rangle = 6\mu kT\,t = 6Dt$ met $D = \mu kT$ (\textbf{Einstein-diffusierelatie}).\\
Bol met straal $r$: $\mu = 1/(6\pi\eta r)$ (Stokes), dus $\langle l_B^2\rangle = \frac{kT}{\pi\eta r}t$.\\
$\Rightarrow$ Meting van $\langle l_B^2\rangle$ geeft $N_A$ (via $k = R/N_A$).
\end{oplossing}

\noindent\textcolor{headerbg}{\rule{\linewidth}{1pt}}

%===========================================================================
\section{Oefeningen (Open Boek) -- Geordend naar Frequentie}
%===========================================================================

%-------------------------------------------------------------------
\subsection{Carnotcyclus: Arbeid, Warmte en Rendement \textnormal{\color{warningred}[$\star\star\star\star\star$\ -- 5+ examens]}}

\begin{oefening}{Carnot cyclus met een ideaal gas (of fotonengas). \textnormal{(2010, 2016, 2022, 2023, 2024)}}
\textbf{Typische vorm:} Een Carnotmachine werkt tussen $T_H$ en $T_K$. Het werkzame gas doorloopt: isotherme expansie $V_1 \to V_2$ bij $T_H$, adiabatische expansie $V_2 \to V_3$, isotherme compressie $V_3 \to V_4$ bij $T_K$, adiabatische compressie $V_4 \to V_1$.\\
(a) Bereken $W$ en $Q$ in elk pad. (b) Toon aan dat $\eta = 1 - T_K/T_H$. (c) Soms: fotonengas of variant.
\end{oefening}

\begin{oplossing}
\textbf{Isotherm bij $T_H$:} $W_{12} = -nRT_H\ln(V_2/V_1) < 0$,\; $Q_H = nRT_H\ln(V_2/V_1) > 0$,\; $\Delta U = 0$.

\textbf{Adiabaat $T_H \to T_K$:} $Q=0$,\; $W_{23} = C_V(T_K-T_H) = \frac{nR}{\gamma-1}(T_K-T_H)$. (Alternatief: $W = \frac{P_3V_3 - P_2V_2}{\gamma-1}$)

\textbf{Isotherm bij $T_K$:} $W_{34} = nRT_K\ln(V_3/V_4)$... maar via adiabaten: $T_HV_2^{\gamma-1} = T_KV_3^{\gamma-1}$ en $T_HV_1^{\gamma-1} = T_KV_4^{\gamma-1}$ $\Rightarrow$ $V_2/V_1 = V_3/V_4$.

Dus: $|Q_K| = nRT_K\ln(V_2/V_1)$ en
$$\eta = \frac{|Q_H| - |Q_K|}{|Q_H|} = \frac{nRT_H\ln(V_2/V_1) - nRT_K\ln(V_2/V_1)}{nRT_H\ln(V_2/V_1)} = 1 - \frac{T_K}{T_H}$$

\textbf{Netto arbeid:} $|W| = (T_H - T_K)nR\ln(V_2/V_1) = (T_H-T_K)(S_2-S_1)$

\textbf{Fotonengas variant (2010-2011):} $P = w/3$, $U = wV$, $PV^{4/3} = \text{const}$ (adiabaat). Gebruik $TdS = C_V\,dT + T\pder{P}{T}_V dV$ met $C_V = 4bVT^3$ en $\pder{P}{T}_V = \frac{4b}{3}T^3$.

\textbf{Variant: monoatomisch vs.\ diatomisch (2022-2023):}\\
Netto arbeid $|W| = |Q_H|\cdot\eta = nRT_H\ln(V_\text{max}/V_\text{min})\cdot(1-T_K/T_H)$.\\
$C_V$ bepaalt de adiabaten ($\gamma$): $V_\text{min}$ en $V_\text{max}$ veranderen, maar $|W| = nRT_H(1-T_K/T_H)\ln(V_2/V_1)$ {\em als}  $V_1, V_2$ vaststaan. Check of volumes of temperaturen vastliggen!
\end{oplossing}

%-------------------------------------------------------------------
\subsection{Koelmachine / Entropie Principe: Minimale Arbeid \textnormal{\color{warningred}[$\star\star\star\star$\ -- 4 examens]}}

\begin{oefening}{Minimale arbeid voor koelen of het entropieprincipe. \textnormal{(2009, set 4, set 5)}}
\textbf{Typisch:} Twee identieke systemen met $C_P$, initieel op $T_i$. Een koelmachine koelt één tot $T_f < T_i$. Druk constant. Geen faseverandering. Bepaal de minimale arbeid $W_\text{min}$.
\end{oefening}

\begin{oplossing}
\textbf{Methode via entropieprincipe:}\\
Entropie-verandering systeem (koud lichaam): $\Delta S_1 = C_P\ln(T_f/T_i) < 0$.\\
Entropie-verandering warmtereservoir (het andere systeem, ook op $T_i$, neemt warmte op bij $T_i$):
$$\text{Warmte opgenomen door warm systeem: } Q_H = Q_K + W \quad (1e \text{ wet})$$
$$Q_K = C_P(T_i - T_f) \quad (\text{warmte onttrokken aan koud systeem})$$
$$\Delta S_2 = \frac{Q_H}{T_i} = \frac{Q_K + W}{T_i}$$
Entropieprincipe: $\Delta S_1 + \Delta S_2 \geq 0$:
$$C_P\ln\frac{T_f}{T_i} + \frac{C_P(T_i-T_f) + W}{T_i} \geq 0$$
$$W \geq C_P\left[T_i - T_f + T_i\ln\frac{T_f}{T_i}\right] = C_P\left[T_i\ln\frac{T_i}{T_f} - (T_i-T_f)\right]$$
$$\boxed{W_\text{min} = C_P\left[T_i\ln\frac{T_i}{T_f} - (T_i - T_f)\right]}$$
\textbf{Check: $W_\text{min} > 0$?} Stel $x = T_f/T_i < 1$: $W_\text{min}/C_PT_i = -\ln x - (1-x) = -\ln x + x - 1 > 0$ voor $x<1$. ✓
\end{oplossing}

%-------------------------------------------------------------------
\subsection{TdS Vergelijkingen: Joule-Kelvin Coëfficiënt \textnormal{\color{warningred}[$\star\star\star$\ -- 3 examens]}}

\begin{oefening}{Bepaal $\mu = (\partial T/\partial P)_H$ via TdS. Speciaal geval: ideaal gas en van der Waals. \textnormal{(2009, 2018, set 5)}}
\end{oefening}

\begin{oplossing}
Vertrek van $dH = TdS + VdP$. Gebruik tweede TdS-vergelijking $TdS = C_P\,dT - T\pder{V}{T}_P dP$:
$$dH = C_P\,dT + \left[V - T\pder{V}{T}_P\right]dP$$
$$\pder{T}{P}_H = \mu = \frac{1}{C_P}\left[T\pder{V}{T}_P - V\right]$$

\textbf{Ideaal gas:} $\pder{V}{T}_P = nR/P = V/T$ $\Rightarrow$ $\mu = \frac{1}{C_P}\left[V - V\right] = 0$. Ideaal gas koelt nooit/warmt nooit bij smoorproces.

\textbf{Van der Waals gas} (1 mol): $P = \frac{RT}{v-b} - \frac{a}{v^2}$:
$$\pder{v}{T}_P = \frac{R}{P + a/v^2 - 2a(v-b)/v^3} \approx \frac{R}{P}\left[1 - \frac{a}{RTv} + \frac{2b}{v} + ...\right]$$
Vereenvoudigd voor lage $P$ (ideaal): $\mu \approx \frac{1}{C_P}\left[\frac{2a}{RT} - b\right]$. 
Inversietemperatuur: $T_\text{inv} = 2a/(Rb)$.
\end{oplossing}

%-------------------------------------------------------------------
\subsection{$C_P - C_V$ voor een Van der Waals Gas \textnormal{\color{warningred}[$\star\star\star$\ -- 3 examens]}}

\begin{oefening}{Bepaal $C_P - C_V$ voor een van der Waals gas in functie van $T$, $v$, $a$ en $b$. \textnormal{(2018, set 2, set 5)}}
\end{oefening}

\begin{oplossing}
Gebruik: $C_P - C_V = -T\pder{V}{T}^2_P \pder{P}{V}^{-1}_T = \frac{TV\beta^2}{\kappa}$

Of direct: $C_P - C_V = T\pder{V}{T}_P\pder{P}{T}_V$ (via Maxwell + partiële afgeleiden).

\textbf{Voor van der Waals} ($P = \frac{RT}{v-b} - \frac{a}{v^2}$):
$$\pder{P}{T}_V = \frac{R}{v-b}, \quad \pder{P}{V}_T = -\frac{RT}{(v-b)^2} + \frac{2a}{v^3}$$
$$\pder{V}{T}_P = -\frac{(\partial P/\partial T)_V}{(\partial P/\partial V)_T} = \frac{R/(v-b)}{RT/(v-b)^2 - 2a/v^3}$$
$$C_P - C_V = T \cdot \frac{R}{v-b} \cdot \pder{V}{T}_P = \frac{TR^2/(v-b)^2}{RT/(v-b)^2 - 2a/v^3} = \frac{R}{1 - \frac{2a(v-b)^2}{RTv^3}}$$

Voor lage $P$ (groot $v$): $C_P - C_V \approx R\left[1 + \frac{2a(v-b)^2}{RTv^3}\right] > R$.\\
In het kritisch punt: divergeert ($2a(v_K-b)^2/(RT_Kv_K^3) = 1$).
\end{oplossing}

%-------------------------------------------------------------------
\subsection{Stefan-Boltzmann / Straling: Toepassingen \textnormal{\color{warningred}[$\star\star\star$\ -- 3 examens]}}

\begin{oefening}{Stralings-toepassingen: zonneconstante, CMB fotonen, evaporatie van helium. \textnormal{(2010, 2015, set 6)}}
\end{oefening}

\begin{oplossing}
\textbf{Temperatuur van de zon (2015-2016):}\\
Zonneconstante $\Phi = 1360\,\text{W/m}^2$. Straal zon $R_\odot = 6{,}96\times10^8\,\text{m}$, afstand $d = 1\,\text{AU} = 1{,}496\times10^{11}\,\text{m}$.
$$\Phi = \sigma T_\odot^4 \cdot \frac{R_\odot^2}{d^2} \Rightarrow T_\odot = \left(\frac{\Phi d^2}{\sigma R_\odot^2}\right)^{1/4} \approx 5780\,\text{K}$$

\textbf{CMB fotonen (2010-2011):} $T = 2{,}725\,\text{K}$.\\
Gemiddelde energie per foton: $\langle h\nu\rangle \approx 2{,}7\,kT \approx 6{,}3\times10^{-4}\,\text{eV}$.\\
Energiedichtheid: $w = bT^4 = aT^4$ met $a = 4\sigma/c$.\\
Aantal fotonen per volume: $n = w/\langle h\nu\rangle \approx 410\,\text{cm}^{-3}$.

\textbf{Helium cryostaat (set 6, opgave 11):}\\
Radiatie van stikstofmantel ($T_2 = 80\,\text{K}$, $\varepsilon = 0{,}25$) op He-vat ($T_1 = 4\,\text{K}$):
$$\Phi = A\varepsilon\sigma(T_2^4 - T_1^4) \approx A\varepsilon\sigma T_2^4$$
Verdampte He per uur: $\dot{m} = \Phi/L$ met $L = 21\,\text{kJ/kg}$.
\end{oplossing}

%-------------------------------------------------------------------
\subsection{Thermische Cycli: PV-diagram, Rendement en TS-diagram \textnormal{\color{warningred}[$\star\star\star$\ -- 3 examens]}}

\begin{oefening}{Gasturbine (Brayton), Otto met van der Waals, speciale cycli. \textnormal{(2015, 2016, 2023, 2024)}}
\end{oefening}

\begin{oplossing}
\textbf{Brayton-cyclus (gasturbine, 2015-2016):} 2 isobaren + 2 adiabaten.
\begin{itemize}
  \item PV-diagram: isobaar expansie (warmteopname $Q_H = C_P(T_3-T_2)$), adiabatische expansie, isobaar compressie ($|Q_K| = C_P(T_4-T_1)$), adiabatische compressie.
  \item Met $T_2/T_1 = (P_H/P_L)^{(\gamma-1)/\gamma} = r_P^{(\gamma-1)/\gamma}$:
  $$\eta_\text{Brayton} = 1 - \frac{T_1}{T_2} = 1 - r_P^{-(\gamma-1)/\gamma}$$
\end{itemize}

\textbf{Cyclus als rechthoek in PV-diagram (2023-2024, 2e zit):}\\
Processen: $(P_1,V_1) \to (P_2,V_1) \to (P_2,V_2) \to (P_1,V_2) \to (P_1,V_1)$.
\begin{itemize}
  \item Netto arbeid: $|W| = (P_2-P_1)(V_2-V_1)$
  \item Warmteopname: isochorisch $(V_1,P_1\to P_2)$ en isobaar $(P_2, V_1\to V_2)$: $Q_H = C_V(T_2-T_1) + C_P\cdot P_2(V_2-V_1)/(nR)$
  \item Rendement: $\eta = |W|/Q_H$
\end{itemize}

\textbf{TS-diagram tekenen:} isobaar = curve met $T\propto S$ (uit $dS=C_P dT/T$); isochoor = steilere curve; adiabaat = verticale lijn; isotherm = horizontale lijn.
\end{oplossing}

%-------------------------------------------------------------------
\subsection{Diffusie / Transport: Druk als Functie van de Tijd \textnormal{\color{warningred}[$\star\star\star$\ -- 2+ examens]}}

\begin{oefening}{Diffusie van gas via kleine opening tussen twee compartimenten. \textnormal{(2018, 2019 schriftelijk)}}
Afgesloten vat verdeeld door diathermische wand met opening (oppervlakte $A$). Links: $N$ moleculen, ideaal gas, $T_0$, $P_0$. Rechts: leeg. Bepaal $P(t)$.
\end{oefening}

\begin{oplossing}
\textbf{Flux door opening:} Aantal moleculen per tijdseenheid door $A$: $\Phi_n = nA\bar{v}/4$ met $\bar{v} = \sqrt{8kT/\pi m}$ en $n = N/V$ de molecuuldichtheid.

\textbf{Netto flux} (links naar rechts - rechts naar links):
$$\frac{dN_L}{dt} = -A\frac{(N_L - N_R)\bar{v}}{4V}$$
Met $N_L + N_R = N$: stel $\Delta N = N_L - N_R$:
$$\frac{d\Delta N}{dt} = -\frac{A\bar{v}}{2V}\Delta N \Rightarrow \Delta N(t) = N\,e^{-t/\tau}, \quad \tau = \frac{2V}{A\bar{v}}$$
$$N_L(t) = \frac{N}{2}(1 + e^{-t/\tau}), \quad N_R(t) = \frac{N}{2}(1 - e^{-t/\tau})$$
$$\boxed{P_L(t) = \frac{N_L kT}{V} = \frac{P_0}{2}(1+e^{-t/\tau}), \quad P_R(t) = \frac{P_0}{2}(1-e^{-t/\tau})}$$
Evenwicht: $P_L = P_R = P_0/2$. De thermische evenwichtstemperatuur $T_0$ blijft constant (diathermische wand).
\end{oplossing}

%-------------------------------------------------------------------
\subsection{Warmtegeleiding: Stationaire Toestand en IJsgroei \textnormal{\color{warningred}[$\star\star\star$\ -- 2 examens]}}

\begin{oefening}{Warmtegeleiding door ijslaag: snelheid van groei. \textnormal{(2015, 2016)}}
Water op $0\,°C$, boven ijslaag (dikte $x$), buitentemperatuur $T_b < 0\,°C$. Bepaal hoe snel de ijslaag groeit.
\end{oefening}

\begin{oplossing}
\textbf{Warmtestroom door ijslaag} (stationaire toestand):
$$\Phi = \lambda_\text{ijs}\, A\, \frac{|T_b|}{x}$$
De warmtestroom is gelijk aan de latente smelt/stollingssnelheid:
$$\Phi = L\,\rho_\text{ijs}\, A\, \frac{dx}{dt}$$
Gelijkstellen:
$$\frac{dx}{dt} = \frac{\lambda_\text{ijs}\,|T_b|}{L\,\rho_\text{ijs}\,x} \Rightarrow x\,dx = \frac{\lambda_\text{ijs}|T_b|}{L\rho_\text{ijs}}dt$$
$$\boxed{x(t) = \sqrt{\frac{2\lambda_\text{ijs}|T_b|}{L\rho_\text{ijs}}\,t}}$$
Typische waarden: $\lambda_\text{ijs} = 2\,\text{W/(m\,K)}$, $L = 333\,\text{kJ/kg}$, $\rho_\text{ijs} = 917\,\text{kg/m}^3$.
\end{oplossing}

%-------------------------------------------------------------------
\subsection{Statistische Mechanica: Temperatuur Schatten en Partitiefunctie \textnormal{\color{warningred}[$\star\star$\ -- 2 examens]}}

\begin{oefening}{Schat de temperatuur uit een energieverdeling of entropiewaarden. \textnormal{(2023, 2024 eerste zit)}}
(a) Gegeven bezettingspercentages bij energieniveaus $\varepsilon_i$. (b) Gegeven $S(t)$ voor een systeem.
\end{oefening}

\begin{oplossing}
\textbf{(a) Via Boltzmann-verdeling:}\\
$N_i/N_j = (g_i/g_j)\,e^{-(\varepsilon_i-\varepsilon_j)/kT}$. Neem twee niveaus:
$$\frac{N_i}{N_j} = e^{-(\varepsilon_i-\varepsilon_j)/kT} \Rightarrow kT = -\frac{\varepsilon_i-\varepsilon_j}{\ln(N_i/N_j)}$$
Met de gegeven waarden: bijv.\ $N_1/N_2 = 3{,}1\%/63\% = 0{,}049$, $\varepsilon_1 - \varepsilon_2 = (30{,}1-4{,}3)\,\text{meV} = 25{,}8\,\text{meV}$:
$$kT = \frac{25{,}8}{-\ln(0{,}049)} = \frac{25{,}8}{3{,}02} \approx 8{,}5\,\text{meV} \Rightarrow T \approx 100\,\text{K}$$

\textbf{(b) Via $1/T = (\partial S/\partial U)_V$:}\\
$S(t)$ gegeven na 500 s: bij evenwicht ($t\to\infty$) nadert $S$ een constante. Bepaal $dS/dt$ en het plateau.

\textbf{Heliumatomen in doos (2023):} $N = 6\times10^{22}$, $V = 2\,\text{l}$, $P = 1\,\text{atm}$.
$T = PV/(Nk) = 101325\times 0{,}002/(6\times10^{22}\times1{,}38\times10^{-23}) \approx 244\,\text{K}$
\end{oplossing}

%-------------------------------------------------------------------
\subsection{Calorimetrie en Smeltwarmte \textnormal{\color{warningred}[$\star\star$\ -- 2 examens]}}

\begin{oefening}{1 kg kwik bij smeltpunt in calorimeter (0,50 kg Al + 1,2 kg water bij 20,0 °C). Evenwicht bij 16,5 °C. Bepaal smeltwarmte van kwik. \textnormal{(2009, set 1)}}
\end{oefening}

\begin{oplossing}
\textbf{Warmtebalans:} warmte gewonnen door kwik = warmte verloren door water + Al.

Warmte verloren door water: $Q_\text{water} = m_w c_w (T_i - T_f) = 1{,}2\times4186\times(20-16{,}5) = 17{,}58\,\text{kJ}$\\
Warmte verloren door Al: $Q_\text{Al} = m_\text{Al}c_\text{Al}(T_i-T_f) = 0{,}5\times900\times3{,}5 = 1{,}575\,\text{kJ}$\\
Warmte opgenomen door kwik (smelt + opwarmen van $-39$ tot $16{,}5\,°C$):
$$Q_\text{Hg} = m_\text{Hg}L_\text{Hg} + m_\text{Hg}c_\text{Hg}(T_f - T_\text{smelt}) = 1\times L + 1\times140\times55{,}5$$
$$Q_\text{water} + Q_\text{Al} = Q_\text{Hg}: \quad 17580 + 1575 = L + 7770$$
$$L = 11{,}385\,\text{kJ/kg} \approx 11{,}4\,\text{kJ/kg}$$
\textit{Literatuurwaarde: 11,8 kJ/kg. }
\end{oplossing}

%-------------------------------------------------------------------
\subsection{Thermodynamische Functies: Gibbsfunctie en Toestandsvergelijking \textnormal{\color{warningred}[$\star\star$\ -- 2 examens]}}

\begin{oefening}{Gegeven is de Gibbsfunctie $G(T,P)$ van een gas. Leid de toestandsvergelijking af. \textnormal{(2022, 2023 gesloten boek)}}
$$G = RT\ln P + A + BP + \frac{CP^2}{2} + \frac{DP^3}{3}$$
\end{oefening}

\begin{oplossing}
Gebruik $V = \pder{G}{P}_T$:
$$V = \pder{G}{P}_T = \frac{RT}{P} + B + CP + DP^2$$
$$\boxed{PV = RT + BP + CP^2 + DP^3}$$
Dit is een viriaalontwikkeling met $A = RT$, $B$, $C$, $D$ als viriaalcoëfficiënten.

\textbf{Bonus:} $S = -\pder{G}{T}_P = -R\ln P - \frac{dA}{dT} - \frac{dB}{dT}P - \ldots$ en $C_P = -T\pder{^2G}{T^2}_P$.
\end{oplossing}

%-------------------------------------------------------------------
\subsection{Zwart Gat Thermodynamica (Recent Populair) \textnormal{\color{warningred}[$\star\star$\ -- recent examen 2024]}}

\begin{oefening}{Bekenstein-Hawking entropie $S = \frac{kc^3A}{4G\hbar}$, $R_S = 2GM/c^2$, $U = Mc^2$. \textnormal{(2023-2024 2e zit)}}
(a) Bepaal $C_V$. (b) Hoelang duurt volledige verdamping?
\end{oefening}

\begin{oplossing}
\textbf{(a) Temperatuur:} $1/T = \pder{S}{U}_V$. Met $U = Mc^2$ en $R_S = 2GM/c^2$:
$$A = 4\pi R_S^2 = \frac{16\pi G^2 M^2}{c^4} = \frac{16\pi G^2 U^2}{c^8}$$
$$S = \frac{kc^3}{4G\hbar}\cdot\frac{16\pi G^2 U^2}{c^8} = \frac{4\pi kG}{c^5\hbar}U^2$$
$$\frac{1}{T} = \pder{S}{U} = \frac{8\pi kG}{c^5\hbar}U = \frac{8\pi kGM}{c^3\hbar}$$
$$T = \frac{\hbar c^3}{8\pi kGM} \quad \text{(Hawking-temperatuur)}$$
$$C_V = \pder{U}{T}_V = c^2\pder{M}{T}_V = -\frac{\hbar c^5}{8\pi kG T^2 c^2} \cdot c^2 = -\frac{c^2}{8\pi kGT^2/(\hbar c^3)} \cdot 1$$

Eenvoudiger: $U = Mc^2$, $T = \frac{\hbar c^3}{8\pi kGM}$ $\Rightarrow$ $M = \frac{\hbar c^3}{8\pi kGT}$, $U = \frac{\hbar c^5}{8\pi kGT}$:
$$C_V = \pder{U}{T} = -\frac{\hbar c^5}{8\pi kG T^2} < 0 \quad \text{(negatief!)}$$

\textbf{(b) Verdampingstijd:} Zwart lichaam straalt $P = -dU/dt = \sigma T^4 \cdot 4\pi R_S^2$.
$$-\frac{dM}{dt} = \frac{\sigma T^4 \cdot 4\pi R_S^2}{c^2} = \frac{4\pi\sigma}{c^2}\left(\frac{\hbar c^3}{8\pi kGM}\right)^4\left(\frac{2GM}{c^2}\right)^2 \propto M^{-2}$$
Integreer: $M^2 dM \propto -dt$ $\Rightarrow$ $t \propto M_0^3$.\\
$t_\text{evap} = \frac{5120\pi G^2 M_0^3}{\hbar c^4}$. Voor $M_0 = 10^{22}$ kg: $t \approx 2\times10^{33}$ jaar.
\end{oplossing}

%-------------------------------------------------------------------
\subsection{Warmtecapaciteit in Gravitatieveld / Energieverdeling \textnormal{\color{warningred}[$\star\star$\ -- 1-2 examens]}}

\begin{oefening}{Moleculen in cilinder, onderhevig aan zwaartekracht: $\varepsilon = \frac{3}{2}kT + mg\bar{z}$. \textnormal{(2022, 2023)}}
Bereken $c_V$ per deeltje. Bespreek $T\to0$ en $T\to\infty$.
\end{oefening}

\begin{oplossing}
\textbf{Partitiefunctie:} Translatie $\times$ zwaartekracht:
$$Z = Z_\text{trans} \cdot Z_\text{grav}, \quad Z_\text{grav} = \int_0^L e^{-mgz/kT}dz = \frac{kT}{mg}(1-e^{-mgL/kT})$$
$$U = NkT^2\pder{\ln Z}{T} = \frac{3}{2}NkT + Nmg\left[\frac{kT}{mg} - \frac{L}{e^{mgL/kT}-1}\right]$$
$$c_V = \pder{U/N}{T} = \frac{3}{2}k + k\left[1 - \frac{(mgL/kT)^2 e^{mgL/kT}}{(e^{mgL/kT}-1)^2}\right]$$

\textbf{Limieten:}
\begin{itemize}
  \item $T\to\infty$: gravitatie verwaarloosbaar, $c_V \to \frac{5}{2}k$ (3 translatie + 1 gravitatie=2 terms, maar limiet: $x\to0$: het tweede term gaat naar 1, dus $c_V\to 5/2\, k$) \quad -- check: equipartitie: 3 translatie + 1 potentieel/kinetisch gravitatie = $\frac{5}{2}k$? Hier $c_V\to\frac{5}{2}k$.
  \item $T\to0$: deeltjes vriezen op bodem, gravitationele vrijheidsgraad ``bevriest'', $c_V \to \frac{3}{2}k$.
\end{itemize}
\end{oplossing}

\noindent\textcolor{headerbg}{\rule{\linewidth}{1pt}}

%===========================================================================
\section{Snelle Referentiegids: Wat Moet Je Echt Kennen?}
%===========================================================================

\begin{multicols}{2}

\begin{tip}
\textbf{THEORIE: Top 5 meest gevraagd}
\begin{enumerate}
  \item Caratheodory + 2e wet (stap voor stap!)
  \item Entropie: 2 definities + $S$ ideaal gas
  \item Carnot + Kelvin $T$-schaal
  \item TdS-vergelijkingen + toepassingen
  \item Planck + Stefan-Boltzmann
\end{enumerate}
\end{tip}

\begin{tip}
\textbf{OEFENINGEN: Top 5 meest gevraagd}
\begin{enumerate}
  \item Carnot: $W$, $Q$, $\eta$ berekenen
  \item Minimale arbeid koelmachine (entropie)
  \item Joule-Kelvin coëfficiënt via TdS
  \item $C_P - C_V$ (vdW of algemeen)
  \item Straling (Stefan/zonneconstante/CMB)
\end{enumerate}
\end{tip}

\begin{formulabox}[title=Handige Identiteiten]
$\pder{x}{y}_z = \left[\pder{y}{x}_z\right]^{-1}$\\[4pt]
$\pder{x}{y}_z\pder{y}{z}_x\pder{z}{x}_y = -1$\\[4pt]
$dF = -SdT - PdV \Rightarrow$ F vrij bij isochoor-isotherm\\[4pt]
$dG = -SdT + VdP \Rightarrow$ G vrij bij isobaar-isotherm
\end{formulabox}

\begin{formulabox}[title=Van der Waals]
$(P+a/v^2)(v-b) = RT$\\
$P_K = a/(27b^2)$, $v_K = 3b$, $T_K = 8a/(27Rb)$\\
$\pder{u}{v}_T = a/v^2$
\end{formulabox}

\begin{formulabox}[title=Molaire Warmtecapaciteiten]
Monoatomisch: $c_V = \frac{3}{2}R$, $c_P = \frac{5}{2}R$, $\gamma = \frac{5}{3}$\\
Diatomisch (rot.): $c_V = \frac{5}{2}R$, $c_P = \frac{7}{2}R$, $\gamma = \frac{7}{5}$\\
Diatomisch (rot.+vibr.): $c_V = \frac{7}{2}R$, $\gamma = \frac{9}{7}$
\end{formulabox}

\begin{formulabox}[title=Entropie Irreversibele Processen]
Isotherme dissipatie: $\Delta S_\text{univ} = W/T$\\
Adiabatische dissipatie: $\Delta S = C_P\ln(T_f/T_i)$\\
Vrije expansie ideaal gas: $\Delta S = nR\ln(V_f/V_i)$\\
Warmte-overdracht: $\Delta S_\text{univ} = Q(1/T_2 - 1/T_1)$
\end{formulabox}

\begin{formulabox}[title=Clapeyron \& Faseover-gangen]
$\frac{dP}{dT} = \frac{L}{T\Delta V}$\\
Smelting water: helling negatief ($\Delta V < 0$)\\
Tripelpunt: 3 lijnen snijden samen\\
2e orde: geen $L$, geen $\Delta V$; Ehrenfest: $\frac{dP}{dT} = \frac{\Delta C_P}{TV\Delta\beta} = \frac{\Delta\beta}{\Delta\kappa}$
\end{formulabox}

\begin{formulabox}[title=Transport -- Snelle Formules]
$D = \frac{1}{3}\bar{v}\ell$ (diffusie)\\
$\lambda = \frac{1}{3}n\bar{v}\ell c'_V$ (warmtegeleiding)\\
$\eta_\text{visc} = \frac{1}{3}nm\bar{v}\ell$ (viscositeit)\\
$\ell = \frac{1}{\sqrt{2}n\sigma}$, $\bar{v} = \sqrt{8kT/\pi m}$\\
Poiseuille: $Q_D = \frac{\pi R^4 (P_1-P_2)}{8\eta L}$
\end{formulabox}

\end{multicols}

%===========================================================================
\section{Examentips}
%===========================================================================
\begin{tip}
\begin{itemize}
  \item \textbf{Gesloten boek:} Oefen het uitschrijven van Caratheodory van voor naar achter, inclusief het UXX'-diagram. Dit is een van de moeilijkste bewijzen en vraagt oefening.
  \item \textbf{Open boek:} Begin altijd met de eerste wet en schrijf duidelijk welke grootheden vastgehouden worden. Gebruik toestandsvergelijking vroeg in de berekening.
  \item \textbf{Entropie:} Bij irreversibele processen: kies altijd een vervangend reversibel pad dat dezelfde begin- en eindtoestand heeft.
  \item \textbf{Carnotcyclus:} Teken altijd het PV-diagram EN het TS-diagram. In het TS-diagram is de cyclus een rechthoek.
  \item \textbf{Dimensies:} Controleer altijd je eenheden. $[kT] = \text{J}$, $[k] = \text{J/K}$, $[R] = \text{J/(mol\,K)}$.
  \item \textbf{TdS-vergelijkingen:} Leer beide formules van buiten. Ze zijn het vertrekpunt voor heel veel afleidingen.
  \item \textbf{Statistische mechanica:} de partitiefunctie $Z$ is alles. Zodra je $Z$ hebt, volgen $U$, $P$, $S$, $F$ direct.
  \item \textbf{Recent examen (2024):} Trend naar meer toepassingen buiten de cursus (zwarte gaten, astrophysica). Pas thermodynamica toe op onbekende systemen.
\end{itemize}
\end{tip}

\end{document}